% !TEX TS-program = XeLaTeX
% use the following command:
% all document files must be coded in UTF-8
\documentclass[portuguese]{textolivre}
% build HTML with: make4ht -e build.lua -c textolivre.cfg -x -u article "fn-in,svg,pic-align"

\journalname{Texto Livre}
\thevolume{19}
%\thenumber{1} % old template
\theyear{2026}
\receiveddate{\DTMdisplaydate{2025}{10}{4}{-1}} % YYYY MM DD
\accepteddate{\DTMdisplaydate{2025}{11}{30}{-1}}
\publisheddate{\DTMdisplaydate{2026}{6}{12}{-1}}
\corrauthor{Tania Mikaela Garcia Roberto}
\articledoi{10.1590/1983-3652.2026.62181}
%\articleid{NNNN} % if the article ID is not the last 5 numbers of its DOI, provide it using \articleid{} commmand 
% list of available sesscions in the journal: articles, dossier, reports, essays, reviews, interviews, editorial
\articlesessionname{articles}
\runningauthor{Roberto} 
%\editorname{Leonardo Araújo} % old template
\sectioneditorname{Daniervelin Pereira~\orcid{0000-0003-1861-3609}}
\layouteditorname{Leonardo Araujo~\orcid{0000-0003-3884-2177}}

\title{Entre algoritmos e símbolos: ortografia, inteligência artificial e \textit{fake news}}
\othertitle{Between algorithms and symbols: orthography, artificial intelligence, and fake news}
% if there is a third language title, add here:
%\othertitle{Artikelvorlage zur Einreichung beim Texto Livre Journal}

\author[1]{Tania Mikaela Garcia Roberto~\orcid{0000-0002-6339-7602}\thanks{Email: \href{mailto:tmgroberto@gmail.com}{tmgroberto@gmail.com}}}
\affil[1]{Universidade Federal Rural do Rio de Janeiro, Seropédica, RJ, Brasil.}

\addbibresource{article.bib}
% use biber instead of bibtex
% $ biber article

% used to create dummy text for the template file
\definecolor{dark-gray}{gray}{0.35} % color used to display dummy texts
\usepackage{lipsum}
\SetLipsumParListSurrounders{\colorlet{oldcolor}{.}\color{dark-gray}}{\color{oldcolor}}

% used here only to provide the XeLaTeX and BibTeX logos
\usepackage{hologo}

% if you use multirows in a table, include the multirow package
\usepackage{multirow}

% provides sidewaysfigure environment
\usepackage{rotating}

% CUSTOM EPIGRAPH - BEGIN 
%%% https://tex.stackexchange.com/questions/193178/specific-epigraph-style
\usepackage{epigraph}
\renewcommand\textflush{flushright}
\makeatletter
\newlength\epitextskip
\pretocmd{\@epitext}{\em}{}{}
\apptocmd{\@epitext}{\em}{}{}
\patchcmd{\epigraph}{\@epitext{#1}\\}{\@epitext{#1}\\[\epitextskip]}{}{}
\makeatother
\setlength\epigraphrule{0pt}
\setlength\epitextskip{0.5ex}
\setlength\epigraphwidth{.7\textwidth}
% CUSTOM EPIGRAPH - END

% to use IPA symbols in unicode add
%\usepackage{fontspec}
%\newfontfamily\ipafont{CMU Serif}
%\newcommand{\ipa}[1]{{\ipafont #1}}
% and in the text you may use the \ipa{...} command passing the symbols in unicode

% LANGUAGE - BEGIN
% ARABIC
% for languages that use special fonts, you must provide the typeface that will be used
% \setotherlanguage{arabic}
% \newfontfamily\arabicfont[Script=Arabic]{Amiri}
% \newfontfamily\arabicfontsf[Script=Arabic]{Amiri}
% \newfontfamily\arabicfonttt[Script=Arabic]{Amiri}
%
% in the article, to add arabic text use: \textlang{arabic}{ ... }
%
% RUSSIAN
% for russian text we also need to define fonts with support for Cyrillic script
% \usepackage{fontspec}
% \setotherlanguage{russian}
% \newfontfamily\cyrillicfont{Times New Roman}
% \newfontfamily\cyrillicfontsf{Times New Roman}[Script=Cyrillic]
% \newfontfamily\cyrillicfonttt{Times New Roman}[Script=Cyrillic]
%
% in the text use \begin{russian} ... \end{russian}
% LANGUAGE - END

% EMOJIS - BEGIN
% to use emoticons in your manuscript
% https://stackoverflow.com/questions/190145/how-to-insert-emoticons-in-latex/57076064
% using font Symbola, which has full support
% the font may be downloaded at:
% https://dn-works.com/ufas/
% add to preamble:
% \newfontfamily\Symbola{Symbola}
% in the text use:
% {\Symbola }
% EMOJIS - END

% LABEL REFERENCE TO DESCRIPTIVE LIST - BEGIN
% reference itens in a descriptive list using their labels instead of numbers
% insert the code below in the preambule:
%\makeatletter
%\let\orgdescriptionlabel\descriptionlabel
%\renewcommand*{\descriptionlabel}[1]{%
%  \let\orglabel\label
%  \let\label\@gobble
%  \phantomsection
%  \edef\@currentlabel{#1\unskip}%
%  \let\label\orglabel
%  \orgdescriptionlabel{#1}%
%}
%\makeatother
%
% in your document, use as illustraded here:
%\begin{description}
%  \item[first\label{itm1}] this is only an example;
%  % ...  add more items
%\end{description}
% LABEL REFERENCE TO DESCRIPTIVE LIST - END


% add line numbers for submission
%\usepackage{lineno}
%\linenumbers

\begin{document}
\maketitle

\begin{polyabstract}
\begin{abstract}
A ortografia, em tempos digitais, ultrapassa a dimensão normativa para se constituir como campo de disputas simbólicas que articula poder, tecnologia e educação. O objetivo deste artigo é discutir a ortografia como instrumento de exclusão e como possibilidade de emancipação, especialmente quando atravessada pela cultura da desinformação, pelas \textit{fake news} e pelos algoritmos que corrigem, antecipam e avaliam a escrita. De caráter teórico-reflexivo, a pesquisa mobiliza aportes de diferentes áreas do conhecimento, avançando para além da tradicional ênfase nas correspondências entre fonemas e grafemas, que marcam grande parte dos estudos sobre o ensino reflexivo da ortografia. As reflexões apontam que a ortografia, ao mesmo tempo em que é usada como marcador de credibilidade e exclusão, pode ser ressignificada pedagogicamente. O erro, tradicionalmente entendido como falha, é recolocado no presente artigo como espaço de autoria, reflexão e resistência frente à automatização algorítmica e à lógica da pós-verdade. Além disso, destaca-se que a correção formal não garante veracidade, mas opera como fator cognitivo de ilusão de verdade, o que demanda uma pedagogia crítica capaz de articular ortografia, cidadania e emancipação. Conclui-se, assim, que ensinar ortografia na era digital exige integrar conhecimento linguístico e reflexão crítica, assumindo o horizonte dos multiletramentos como campo formativo na área. O artigo sugere que futuras pesquisas explorem estratégias pedagógicas para implementação dessas propostas em contextos concretos, como o do Profletras, aproximando teoria, prática docente e recepção crítica da escrita digital.

\keywords{Ortografia \sep Erro \sep Algoritmo \sep Fake news \sep Multiletramentos}
\end{abstract}

\begin{english}
\begin{abstract}
In digital times, orthography transcends its normative dimension to become a field of symbolic disputes that articulate power, technology, and education. The aim of this article is to discuss orthography both as an instrument of exclusion and as a possibility for emancipation, especially when intertwined with the culture of disinformation, fake news, and the algorithms that correct, anticipate, and evaluate writing. With a theoretical-reflective character, the research draws upon contributions from different areas of knowledge, moving beyond the traditional emphasis on phoneme-grapheme correspondences that have shaped much of the studies on reflective orthography teaching. The reflections suggest that orthography, while functioning as a marker of credibility and exclusion, can also be pedagogically re-signified. Error, traditionally understood as failure, is reframed here as a space of authorship, reflection, and resistance against algorithmic automation and the logic of post-truth. Furthermore, the study highlights that formal correction does not guarantee truth but operates as a cognitive factor of illusion of truth, which requires a critical pedagogy capable of articulating orthography, citizenship, and emancipation. It concludes that teaching orthography in the digital age demands integrating linguistic knowledge with critical reflection, assuming multiliteracies as a formative horizon. The article also suggests that future research should explore pedagogical strategies for implementing these proposals in concrete contexts, such as Profletras, bridging theory, teaching practice, and critical reception of digital writing.

\keywords{Orthography \sep Error \sep Algorithm \sep Fake news \sep Multiliteracies}
\end{abstract}
\end{english}
% if there is another abstract, insert it here using the same scheme
\end{polyabstract}

\section{Introdução}\label{sec-intro}
Nos gestos cotidianos de escrever uma mensagem rápida no celular, revisar um e-mail profissional ou postar em uma rede social, a ortografia está sempre presente, silenciosa, mas decisiva. Pequenos deslizes podem ser lidos como descuido; correções automáticas, aceitas sem reflexão; e textos impecáveis são, quase sempre, tomados como mais confiáveis. A escrita, que parece simples, nunca foi neutra: ela carrega expectativas, julgamentos e, como bem sabemos, preconceitos.

Essa constatação, vale destacar, não é nova. Bakhtin \cite{volochinov2017marxismo}\footnote{Valentin Volóchinov é um pseudônimo de Mikhail Bakhtin usado na primeira edição brasileira da obra \textit{Marxismo e a filosofia da linguagem} traduzida da publicação russa de 1929.} já lembrava que todo ato de linguagem é atravessado por valores e ideologias, impossibilitando qualquer tipo de neutralidade. \textcite{fairclough1989language,fairclough2001discurso}, por sua vez, destacou que os textos são sempre práticas sociais constitutivas de relações de poder, e não meras representações inocentes da realidade. No contexto brasileiro, \textcite{bagno1999preconceito,bortoni2004linguamaterna,bortoni2005chegemu} reforçam essa perspectiva quando evidenciam como a valorização da norma culta e a estigmatização das variedades linguísticas funcionam como instrumentos de exclusão e (des)legitimação social. Assim, longe de ser apenas um código gráfico, a escrita é também um campo de disputa simbólica.

É nesse cruzamento que se encontram algoritmos e símbolos. Os algoritmos – corretores, teclados preditivos ou sistemas de IA – atuam como mediadores permanentes da escrita, frequentemente corrigindo erros antes mesmo que o usuário tenha a chance de refletir sobre eles. \textcite{gillespie2014algorithms} mostra como algoritmos, ao definir o que se torna visível ou invisível, participam ativamente da construção da experiência social. E, nessa mesma direção, \textcite{noble2018algorithms} mostra que esses processos estão longe da neutralidade: ao organizar informações, os algoritmos podem reforçar desigualdades históricas e transformar preconceitos sociais em padrões digitais. Nesse cenário, o erro ortográfico, frequentemente apagado pelas correções automáticas, deixa de ser apenas uma falha individual e passa a se configurar como terreno decisivo de disputa pedagógica e social, aspecto que será desenvolvido ao longo deste artigo. Ainda que a análise de \textcite{noble2018algorithms} não trate especificamente da ortografia, sua crítica evidencia como algoritmos reproduzem desigualdades, o que pode ser analogamente pensado no campo da escrita.

Já os símbolos – letras, acentos e convenções – estruturam a ortografia, mas também os significados sociais que lhes são atribuídos. Esses símbolos não apenas representam a língua; eles também funcionam como marcadores de distinção social, indicando quem é reconhecido como “legítimo” e quem é desvalorizado por formas estigmatizadas de escrever. Para \textcite{bourdieu1989poder,bourdieu1991language}, o domínio da chamada “língua legítima” não significa apenas seguir regras, mas ocupar uma posição de legitimidade dentro de jogos sociais e políticos que naturalizam hierarquias e exclusões.

Esses conceitos, embora inicialmente complexos, serão retomados ao longo do artigo com exemplos concretos e implicações pedagógicas, especialmente no que diz respeito ao papel do erro nos multiletramentos, entendidos como práticas que envolvem múltiplas linguagens, mídias e contextos culturais \cite{newlondongroup1996multiliteracies,rojo2009letramentos,rojo2012multiletramentos}. Essa retomada busca tornar clara a pertinência desses conceitos no debate sobre ortografia em tempos digitais. Até porque, na contemporaneidade, os desafios ligados ao ensino da ortografia não se restringem à tensão entre norma e variação.

Vale ressaltar, também, que vivemos sob a lógica da pós-verdade, em que a ortografia assume novos papéis: textos impecáveis podem legitimar desinformações, como mencionam \textcite{fazio2015illusorytruth}, enquanto deslizes mínimos podem ser explorados como armas de ridicularização e estigmatização. É nesse sentido que, entre a automatização algorítmica, que redefine o lugar do erro, e a manipulação discursiva, que explora a forma como critério de credibilidade, a ortografia deixa de ser uma formalidade e passa a ter relevância estratégica, tanto política quanto pedagógica.

Este artigo tem como objetivo, portanto, discutir a ortografia em sua dupla condição: como instrumento de exclusão e como possibilidade de emancipação, sobretudo quando pensada no cruzamento entre tecnologias digitais e a chamada cultura da desinformação. A justificativa para esse debate consiste na necessidade de repensar o ensino da ortografia em um contexto em que a escrita se massifica nas redes e, ao mesmo tempo, pode ser usada como marcador de credibilidade. Ignorar esse quadro é abrir mão da chance de formar sujeitos capazes de compreender a dimensão social da norma e de atuar criticamente em uma sociedade saturada de informação.

É importante lembrar, contudo, que o debate sobre o ensino da ortografia no Brasil já conta com um campo consolidado de pesquisas sobre o ensino de ortografia. Trago aqui três nomes representativos, ainda que não exaustivos, de pesquisadores relevantes, sem contar os que retomam os estudos sobre o ensino de ortografia desde o início das atividades do Profletras\footnote{O Mestrado Profissional em Letras (Profletras) é um programa de pós-graduação em rede nacional, coordenado pela Universidade Federal do Rio Grande do Norte (UFRN) em parceria com a CAPES, voltado para a formação continuada de professores de Língua Portuguesa do ensino fundamental, com ênfase na articulação entre pesquisa e prática pedagógica e que é oferecido por 58 unidades em todo o território nacional \cite{ufrn2025profletras}} em 2013. Maria Bernadete Abaurre, linguista com destacadas pesquisas em fonologia e aquisição da escrita, dedicou-se a investigar a fase inicial da alfabetização, especialmente como as crianças constroem hipóteses sobre o funcionamento do sistema de escrita \cite{abaurre1986alfabetizacao,abaurre1991segmentacao,abaurre1999refaccao,abaurre2011escritaespontanea,abaurrecoudry2008afasicos,abaurre1997cenas}, trazendo relevantes contribuições para os estudos de ortografia. Já \textcite{scliarcabral2003a_guia,scliarcabral2003b_principios,scliarcabral2011desafios,scliarcabral2015melhoria,scliarcabral2019politicas,scliarcabral2022falhado}, em trajetória voltada tanto à fundamentação teórica quanto à intervenção prática, tem se dedicado a discutir a alfabetização a partir dos princípios do sistema alfabético do português, articulando-os às práticas pedagógicas, à consciência fonológica e às políticas públicas, com o objetivo de qualificar o ensino da leitura e da escrita no Brasil. \textcite{morais2002ortografia,morais2005aprendizado}, por sua vez, ocupa papel central nesse contexto ao defender um ensino reflexivo da ortografia, perspectiva que já se fazia presente desde suas contribuições nos Parâmetros Curriculares Nacionais \cite{brasil1997pcnlp}.

Em suas especificidades, esses e outros autores mostram o quanto o ensino da ortografia pode ser pensado para além de uma perspectiva tradicionalista, que a reduz a critério avaliativo e prática de exercícios de fixação, apontando caminhos que a situam como parte do desenvolvimento da consciência linguística e da aprendizagem significativa da língua escrita. Nesse mesmo horizonte, tenho buscado contribuir com pesquisas que articulam o ensino da ortografia aos subsídios da fonologia \cite{roberto2016ortografia,roberto2023caminhos}, inclusive orientando trabalhos acadêmicos que exploram essa interface, como \textcite{cesar2017acrescimo,ramos2018sonoridade,marins2024monotongacao,marinsroberto2025parlaventura}, entre outros.

Neste artigo, contudo, proponho pensar o ensino da ortografia não apenas sob uma ótica reflexiva, ancorada na consciência fonológica e em estratégias metacognitivas. Defendo também uma perspectiva crítica e ampliada que, sem abandonar a dimensão linguística inquestionavelmente relevante, busca municiar o professor a desenvolver um olhar que ultrapasse os limites do nível fonológico, dos princípios do sistema de escrita e das normas convencionais. O que proponho é tomar a ortografia como ponto de partida para refletir sobre relações de poder, preconceito linguístico, políticas linguísticas e processos sociais mais amplos. Tal perspectiva se insere, em alguma medida, nos estudos sobre os multiletramentos, que reconhecem a diversidade de linguagens, mídias e suportes como dimensão constitutiva da cidadania na sociedade digital. O convite é tratar a ortografia não apenas como domínio técnico, mas como campo fértil para explorar sentidos, questionar hierarquias e formar sujeitos críticos diante das práticas linguísticas contemporâneas.

Para desenvolver essa reflexão, dessa forma, o texto se organiza em três movimentos assim articulados: a) a análise do impacto dos algoritmos e das tecnologias digitais na aprendizagem e no uso da ortografia; b) a discussão sobre o papel da ortografia na construção de credibilidade textual e em estratégias de manipulação em tempos de \textit{fake news}; c) a proposição de caminhos para o ensino crítico da ortografia, capaz de integrar norma, tecnologia e práticas de linguagem voltadas a uma formação social e crítica.

\section{Ortografia e tecnologias digitais: entre a automatização e a reflexão crítica}\label{sec-normas}
Seria reducionista pensar a ortografia hoje sem considerar o impacto das tecnologias digitais no processo de aprender e usar a escrita. Corretores automáticos, inteligência artificial e novos letramentos reorganizam práticas de leitura e escrita, alterando, inclusive, as formas de reflexão sobre o erro. Esses dispositivos são mais do que ferramentas técnicas. Eles redesenham as condições de aprendizagem e de circulação da escrita, bem como afetam os critérios sociais de prestígio e credibilidade. \textcite{kalantzis2005learning} já assinalavam que os chamados “novos aprendizados” se constroem na articulação entre diversidade cultural, redes digitais e transformações nos modos de produção do conhecimento, apontando que a educação precisa lidar com múltiplas linguagens, mídias e contextos de uso. Nesse horizonte, a ortografia não pode ser compreendida apenas como norma estática, mas como prática social constantemente tensionada por tecnologias que tornam a escrita mais fluida, automatizada e interativa.

Como observa \textcite{mills2009newlearning},
\textcite{kalantzis2005learning} reconhecem desafios contemporâneos como a globalização, a diversidade cultural e as tecnologias digitais. Esses desafios exigem uma redefinição da educação como ciência. Isso significa concebê-la para além da transmissão de conteúdos, integrando fundamentos teóricos, metodológicos e práticos. Se esse discurso não é novo, se faz mais urgente que ele se consolide na prática pedagógica. Tal perspectiva amplia o alcance da reflexão sobre ortografia, permitindo situá-la no centro de debates que envolvem identidade, inclusão e participação crítica em uma sociedade marcada pela informação em rede.

Como já apontam esses autores, eu diria que ignorar o impacto dessas transformações equivale a permanecer imóvel diante de um \textit{tsunami}: a onda já está em movimento, e ignorar essa onda é se deixar engolir, sem ter construído meios de lidar com ela. O desafio não é negar o novo, mas atravessá-lo criticamente, mergulhando na onda, reconhecendo seus riscos e suas potencialidades. E, nesse cenário, os algoritmos ocupam papel decisivo e precisam ser trazidos ao debate sobre o ensino.

Em termos técnicos, de acordo com \textcite{gillespie2014algorithms}, algoritmos são sequências de instruções matemáticas que orientam máquinas a tomarem decisões e resolver problemas. No cotidiano da escrita, os algoritmos operam de modo invisível em vários níveis. Nos corretores ortográficos, os algoritmos fazem a varredura do texto para identificar combinações de letras que não correspondem ao dicionário interno ou às regras da língua. Em seguida, propõem substituições automáticas que passam a definir o que deve ou não ser considerado “correto”. Nos chamados teclados preditivos – aqueles que sugerem a palavra seguinte enquanto o usuário digita em celulares ou computadores –, os algoritmos analisam padrões estatísticos de uso, tanto do próprio sujeito quanto de grandes bancos de dados, antecipando a forma considerada mais provável antes mesmo que ela seja digitada por completo. Já nos sistemas de inteligência artificial, a operação é ainda mais sofisticada: modelos treinados em amplos \textit{corpora} de textos não apenas sugerem palavras isoladas, mas produzem sequências inteiras, ajustando ortografia, estilo e coerência global do enunciado. Em todos esses casos, a ação algorítmica não se limita a auxiliar a escrita; ela redefine, de forma silenciosa, o que se apresenta como legítimo e aceitável, redefinindo hábitos linguísticos e deslocando a reflexão sobre o erro para uma esfera cada vez mais automatizada.

Essas operações, no entanto, não são infalíveis. Os corretores ortográficos, por exemplo, enfrentam dificuldades quando há mais de uma possibilidade de registro aceitável ou quando o contexto é determinante para o sentido, levando às correções inadequadas que passam despercebidas, já que operam com dicionários fixos e não possuem plena capacidade de interpretar a situação de uso. Os teclados preditivos, por sua vez, que antecipam a palavra provável a partir de padrões estatísticos, nem sempre oferecem resultados pertinentes: não é raro, aliás, que gerem constrangimentos em conversas rápidas ou que se tornem alvo de brincadeiras em redes sociais, nas quais usuários se divertem completando frases apenas com as primeiras sugestões do teclado, evidenciando os limites do sistema. No caso dos modelos de inteligência artificial, embora mais sofisticados, a falibilidade também se faz presente, porque – ainda que os problemas mais acentuados estejam no campo sintático, pragmático e textual –, além de não eliminarem por completo erros ortográficos, podem sugerir formas inadequadas em determinados contextos.

Os algoritmos, assim, estão longe de serem neutros ou perfeitos. Se é verdade que esses sistemas apoiam a escrita, também é verdade que reduzem o espaço de reflexão. Em alguns momentos, a automatização apaga a experiência do erro. Em outros, substituições inadequadas distorcem a reflexão, levando o sujeito a acreditar que errou quando não errou. E mais, como aponta \textcite{noble2018algorithms}, os algoritmos carregam valores e escolhas embutidas em sua programação, e por isso moldam práticas culturais tanto quanto refletem aspectos técnicos, interferindo diretamente em quem tem acesso à visibilidade ou ao apagamento.

A análise de \textcite{noble2018algorithms} mostra, ainda, que motores de busca e sistemas de ordenamento da informação, ao invés de apenas refletirem a realidade social, reproduzem e amplificam estereótipos raciais e de gênero, reforçando desigualdades históricas. Nesse sentido, a ilusão de neutralidade algorítmica esconde mecanismos de poder: o que aparece como resultado confiável é, na verdade, filtrado por lógicas comerciais, políticas e culturais. Assim, ao determinar quais vozes circulam e quais permanecem invisíveis, os algoritmos não apenas organizam dados, mas também participam da construção das hierarquias sociais e simbólicas em que a escrita se insere, numa lógica capitalista e excludente.

A falsa aparência de neutralidade desses sistemas tem se tornado ainda mais evidente em pesquisas recentes sobre rotulagem algorítmica. Um estudo desenvolvido por \textcite{alrawi2022autocomplete} sobre o Google Autocomplete demonstrou que o sistema atribui subtítulos automáticos sistematicamente neutros ou positivos a conspiracionistas, extremistas e até autores de atentados, classificando-os como \textit{writer}, \textit{activist} ou \textit{performer}, apesar dos danos amplamente reconhecidos que suas ações produzem. Embora esses rótulos pareçam tecnicamente corretos, funcionam como filtros simbólicos que suavizam práticas nocivas e reorganizam a percepção pública sobre tais figuras. Esse mecanismo confirma que a ação algorítmica não se limita à escrita, mas constitui formas de mediação social que selecionam o que deve ser visível, crível ou legitimado, reforçando a necessidade de problematizar pedagogicamente aquilo que os sistemas decidem apresentar como “neutro”.

Esse funcionamento desloca o próprio lugar da ortografia. Ao corrigirem automaticamente deslizes ou anteciparem a escrita, os algoritmos reduzem o espaço de enfrentamento com o erro, transformando-o de experiência de reflexão e aprendizagem em operação invisível. Se, em abordagens pedagógicas de ensino reflexivo, nos moldes que propõe \textcite{morais2002ortografia,morais2005aprendizado}, o erro ortográfico era motor de hipóteses e reelaborações, nas plataformas digitais tende a ser suprimido antes mesmo de se consolidar.

\textcite{fernandezjulia2024normadigital} alertam que essa automatização pode neutralizar oportunidades de desenvolvimento, já que aprender envolve lidar com estranhamentos, hesitações e reconstruções, dimensões que a lógica algorítmica tende a suprimir. Para os autores, a chamada “norma digital”, entendida como a padronização promovida por sistemas automáticos de correção e predição, que estabelecem regularidades artificiais no uso da escrita, reduz a possibilidade de o estudante experimentar e refletir sobre suas próprias dificuldades, transformando a aprendizagem em mera adaptação às sugestões do sistema.

Esse processo ameaça empobrecer o espaço de elaboração cognitiva e metalinguística, deslocando o erro de sua função formativa para um plano de invisibilidade ou de correção acrítica. Tal lógica reflete também em uma tendência cultural mais ampla: evitar o erro, o feio, o incômodo. \textcite{vigotski1997formacao} já defendia que o avanço cognitivo ocorre na zona de desenvolvimento proximal, que pressupõe esforço e mediação. Tradicionalmente, abordagens humanistas e dialógicas em educação ressaltam que experiências desafiadoras favorecem processos de crescimento e reorganização interna, articulando-se ao que \textcite{dweck2017mindset} demonstrou ao mostrar que dificuldades podem se converter em oportunidades de aprendizagem. A neurociência reforça esse ponto: estudos recentes mostram que o erro e o desconforto ativam mecanismos cerebrais de monitoramento e ajuste, fundamentais para o desenvolvimento cognitivo \cite{kapur2016productive,ullsperger2014performancemonitoring}. Ao suprimir esses momentos de hesitação, portanto, a lógica algorítmica da “norma digital” não só uniformiza a escrita, mas também nos priva do próprio processo que favorece o aprender.

Pensar caminhos pedagógicos de enfrentamento desse problema abre a possibilidade de transformar essa tensão em recurso, ensinando o aluno a não aceitar de forma passiva as correções automáticas, mas a analisá-las criticamente como parte de sua aprendizagem, caminho a ser discutido mais adiante, neste artigo.

Esse deslocamento, porém, não afeta apenas o processo individual de aprendizagem: ele também repercute no modo como a escrita circula socialmente. Textos sem erros deixam de ser apenas índice de maior escolarização para se tornarem também produto da mediação algorítmica que “normaliza” a escrita. Isso intensifica o valor simbólico da ortografia como marcador de competência, isto é, como sinal socialmente reconhecido de escolarização, inteligência e credibilidade.

Nessa lógica, quem escreve respeitando a norma ortográfica é frequentemente percebido como mais inteligente ou mais apto, independentemente de outras dimensões do discurso. Retomando \textcite{bourdieu1991language}, podemos afirmar que a ortografia funciona como forma de capital linguístico, capaz de conferir prestígio e distinguir grupos sociais. Esse conceito designa o valor social atribuído a determinados usos da língua, funcionando como recurso que pode ser convertido em prestígio e poder simbólico. E, ao ser transferido para a esfera digital, esse capital passa a depender não apenas do domínio escolar, mas também do acesso a dispositivos capazes de garantir correção invisível.

Em outras palavras, escrever bem hoje pode significar duas coisas: domínio escolar e acesso às tecnologias de escrita automatizada. Como versões mais sofisticadas dessas ferramentas são pagas, a imagem de competência linguística passa também pela capacidade de consumir serviços tecnológicos. Isso reforça desigualdades já existentes. Em diálogo com os debates sobre multiletramentos digitais no Brasil \cite{coscarelli2011letramento,rojo2009letramentos,rojo2012multiletramentos,xavier2005letramento}, percebe-se que o domínio da escrita hoje não se restringe a saber lidar com diferentes linguagens e mídias, mas envolve também a inserção desigual em um mercado de ferramentas digitais regulado por interesses capitalistas. Nesse cenário, a ortografia não só mantém seu valor como capital simbólico, mas também se torna índice de acesso diferenciado às condições materiais e tecnológicas de participação social. Todas essas tensões, marcadas por desigualdades e interesses de mercado, podem levar a uma leitura exclusivamente negativa das tecnologias. É importante lembrar, contudo, que os receios diante de inovações não são novos na história, quase toda transformação tecnológica foi inicialmente vista como ameaça antes de ser incorporada de forma crítica e produtiva na vida cotidiana.

A escrita, a imprensa, a eletricidade, o telefone, a televisão, o computador e, mais recentemente, a internet, foram todos recebidos com receios semelhantes em seus momentos históricos. A entrada da calculadora na escola, por exemplo, foi por muito tempo vista como ameaça ao raciocínio matemático. Hoje, entretanto, pesquisas demonstram, como a de \textcite{cunha2018calculadora}, que, quando bem orientada, a calculadora pode potencializar a aprendizagem, estimulando estratégias de resolução de problemas e interpretações críticas de resultados. O paralelo é esclarecedor: tecnologias não anulam o pensamento humano, mas deslocam seus modos de operar.

O desafio, portanto, não é fugir da onda algorítmica que já está em movimento, mas aprender a mergulhar nela de modo crítico, reconhecendo que os algoritmos, embora carreguem riscos e desigualdades, também podem ser convertidos em oportunidade pedagógica e social. É essa tensão – entre o risco de reforço das exclusões e a possibilidade de emancipação – que orienta a análise das próximas seções: primeiro, sobre como a ortografia opera nas estratégias de credibilidade em tempos de \textit{fake news}, e, depois, sobre como podemos pensar caminhos pedagógicos que ressignifiquem o lugar do erro e da norma em uma sociedade marcada pela mediação digital.

\section{Ortografia, \textit{fake news} e a lógica da pós-verdade}\label{sec-conduta}

\epigraph{``Às vezes, a única coisa verdadeira num jornal é a data.''}{Luis Fernando Veríssimo}

Veríssimo expõe, com sua ironia característica e contundente, na epígrafe desta seção, um cenário em que a correção formal de um texto pode criar uma aparência de credibilidade, mesmo quando o conteúdo é manipulado ou falso. A ironia de Veríssimo introduz a problemática da credibilidade textual na pós-verdade. Se, como já explanado na introdução deste artigo, a ortografia é atravessada por valores ideológicos e atua como marcador de legitimidade social \textcite{volochinov2017marxismo,bourdieu1989poder,bourdieu1991language,fairclough1989language,fairclough2001discurso}, seu papel se torna ainda mais visível em tempos de pós-verdade. \textcite{tesich1992government}, citado no editorial de \textcite{siebert2020posverdade}, descreve a pós-verdade como uma condição cultural em que crenças e afetos sobrepõem-se aos fatos, criando terreno fértil para a circulação de desinformações estrategicamente produzidas para moldar percepções sociais. Assim, em meio à lógica da pós-verdade, a ortografia pode ser tanto filtro de exclusão quanto selo de legitimidade, deslocando-se do plano técnico para o centro das disputas simbólicas contemporâneas.

Como enfatiza \textcite{souza2019fakenews}, \textit{fake news} não se reduzem a informações falsas, são discursos planejados para mobilizar emoções, manipular consensos e sustentar projetos de poder. Nesse quadro, deslizes formais são explorados como indícios de incompetência, funcionando como arma de ridicularização e reforço de preconceitos linguísticos. A correção formal atua, assim, como validação, ou seja, textos impecáveis transmitem confiabilidade, mesmo quando carregam distorções ou mentiras.

A psicologia cognitiva ajuda a explicar esse mecanismo. O chamado “efeito da verdade ilusória” \cite{fazio2015illusorytruth} mostra que a repetição aumenta a sensação de veracidade, mesmo quando a informação é falsa. Isso ocorre porque a familiaridade facilita o processamento cognitivo. Assim, quanto mais facilmente algo é lido ou lembrado, mais verdadeiro tende a parecer. A esse quadro se soma o ambiente das chamadas bolhas informacionais \cite{silva2025bolhas}, em que mensagens semelhantes circulam repetidamente em espaços fechados, reforçando crenças pré-existentes e excluindo vozes divergentes.

Nessas condições, a ortografia desempenha um papel decisivo. Pequenos erros formais podem ser usados como indícios de falsidade, desqualificando mensagens ou seus autores, enquanto textos impecáveis, ainda que repletos de distorções, conquistam adesão social. A fluência formal aumenta a facilidade de processamento e, por consequência, a percepção de veracidade \cite{fazio2015illusorytruth}. Isso significa que a ortografia, ao tornar o texto mais fluido e familiar, não apenas sinaliza competência linguística, mas também atua como fator cognitivo que reforça a credibilidade, evidenciando sua centralidade na lógica da desinformação contemporânea.

Esse fenômeno se articula com outros processos de legitimação baseados em fluência formal. A atribuição de subtítulos neutros no Google – como demonstrado no estudo de \textcite{alrawi2022autocomplete} – opera logicamente como a ortografia impecável que mascara conteúdos manipuladores: ambos produzem fluência, facilitam o processamento cognitivo e, por consequência, aumentam a sensação de veracidade. Assim, ao designar conspiracionistas, supremacistas e propagadores de desinformação como “autores”, “apresentadores” ou “empreendedores”, o algoritmo reforça a aparência de normalidade e reorganiza simbolicamente a relação entre forma e credibilidade. Se textos bem escritos podem legitimar mentiras, também rótulos “limpos” podem neutralizar danos sociais graves, evidenciando que, no ecossistema digital, forma e função se combinam para modular crenças e adesões políticas.

Cabe, portanto, à escola trabalhar a ortografia como ferramenta de emancipação, formando sujeitos capazes de desconfiar da correção impecável, interrogar repetições manipuladoras e reconstruir, com base em evidências, os alicerces da vida democrática. O papel da ortografia na pós-verdade confirma seu caráter de símbolo político e pedagógico.


\section{Ensinar ortografia na era digital: entre erros, algoritmos e cidadania crítica}\label{sec-fmt-manuscrito}
Diante do percurso traçado neste artigo – do impacto das tecnologias digitais à relação da ortografia com credibilidade e manipulação discursiva –, a pergunta que devemos nos fazer é: como ensinar ortografia nesse novo cenário? Não se trata mais apenas de transmitir a norma, mesmo que de forma reflexiva, ou focar apenas nas correspondências entre fonemas e grafemas, mas formar sujeitos capazes de compreender os usos sociais da escrita e resistir às formas de exclusão e manipulação que a cercam.

Nesse sentido, não estamos inventando nada de novo, tratamos do campo já consolidado dos letramentos \cite{soares1986letramento,street1995social} e, mais especificamente, dos multiletramentos, tal como formulados pelo \textcite{newlondongroup1996multiliteracies} e aprofundados por autores como \textcite{kalantzis2005learning} e, no Brasil, \textcite{rojo2009letramentos,rojo2012multiletramentos}. O que se propõe aqui é deslocar também a ortografia, tradicionalmente vista como domínio técnico, avaliado por meio de exercícios de fixação, para dentro desse horizonte. Se posso assim chamar, pensar os “letramentos ortográficos” como dimensão formativa que precisa estar presente tanto na alfabetização quanto em níveis posteriores de escolarização, na formação crítica de leitores e produtores de textos em uma sociedade digital marcada por algoritmos, pós-verdade e \textit{fake news}.

Hoje, algoritmos operam invisivelmente e nos apresentam soluções aparentemente perfeitas \cite{gillespie2014algorithms} para a produção textual. Como já exposto, entretanto, esse funcionamento pode neutralizar a experiência cognitiva do erro: ao evitar que ele apareça, suprime-se a oportunidade de refletir sobre o processo de escrita, de formular hipóteses e de desenvolver consciência metalinguística. Em termos freireanos, o erro deixa de ser “forma provisória do saber” \cite{freire1987pedagogia} para tornar-se ruído eliminado pelo sistema, um deslocamento que ameaça a dimensão crítica e, por que não dizer, emancipatória do aprender.

Três aspectos se tornam, portanto, centrais. O primeiro é ressignificar o erro. Como mostram \textcite{ferreiro1985psicogenese}, o erro é parte constitutiva do processo de aprendizagem da escrita, funcionando como hipótese temporária que orienta a construção de saberes, o que se coaduna com o que \textcite{freire1987pedagogia} defende, em uma perspectiva político-pedagógica, ao tratá-lo como “forma provisória do saber”. De modo complementar, \textcite{delatorre2007aprender} propõe uma abordagem multidimensional do erro, entendendo-o como estratégia didática capaz de provocar reflexão, indagações e aprendizagens criativas.

Em um ambiente dominado por algoritmos que corrigem automaticamente, o erro, entretanto, como vimos, é apagado antes de se tornar objeto de análise. Cabe à escola, assim, explorar essa potência pedagógica, propondo situações em que a produção do aluno possa ser colocada em contraste com a “solução” oferecida por corretores automáticos ou pela IA generativa. Tal exercício permite problematizar os vieses desses sistemas e refletir criticamente sobre os limites da automação no campo educacional, em sintonia com a crítica de \textcite{madaio2021beyondfairness} de que a simples avaliação de desempenho não dá conta das injustiças estruturais reproduzidas pela IA.

O segundo aspecto refere-se à massificação da escrita digital. Nunca se produziu tanto texto quanto hoje, em redes sociais, aplicativos de mensagens e plataformas de autopublicação. Essa escrita cotidiana, marcada por flexibilizações normativas e hibridismos entre oralidade e escrita, como apontam \textcite{barton2013language}, não deve ser vista como ameaça, mas pode ser ressignificada como oportunidade pedagógica para repensar o ensino de língua na contemporaneidade. Memes, legendas automáticas e \textit{posts} oferecem material fértil para discutir implicações comunicativas, preconceito linguístico e efeitos sociais da variação ortográfica. Nesse processo, a ortografia pode se tornar ferramenta crítica de leitura do mundo.

Nessa direção, a massificação da escrita digital pode ser incorporada ao ensino não como um desvio a ser corrigido, mas como uma entrada privilegiada para discutir os modos contemporâneos de significação. Do ponto de vista dos multiletramentos, é justamente a circulação em múltiplos gêneros, suportes e registros que permite ao aluno compreender que a ortografia se distribui em diferentes graus de formalidade e expectativa social \cite{newlondongroup1996multiliteracies,rojo2012multiletramentos}. Ao analisar comentários de redes sociais, trechos de fóruns, legendas automáticas ou mensagens de WhatsApp, o professor pode evidenciar que a variação ortográfica não é erro “por falta de regra”, mas adequação a propósitos comunicativos distintos.

Essa abordagem desloca a ortografia da posição de marcador de “certo e errado” para um instrumento de leitura crítica da prática social: por que algumas grafias são interpretadas como descontraídas e outras como incompetentes? Quem define esses critérios? E que efeitos de exclusão se produzem quando a escola desconsidera a diversidade gráfica dos espaços digitais? Trabalhar com esses dados permite aproximar o ensino da ortografia do campo das práticas reais de linguagem, fortalecendo um letramento crítico capaz de desvelar hierarquias simbólicas e preconceitos linguísticos ainda fortemente enraizados nas plataformas digitais.

O terceiro aspecto, articulado aos anteriores, diz respeito à relação entre ortografia, fluência e credibilidade, isto é, à disputa política pelo que conta como verdade em ambientes digitalmente mediados. A mesma escrita digital, que amplia repertórios e flexibiliza usos normativos, também intensifica uma estética de fluência que se converte em índice de credibilidade. Em ambientes saturados por textos curtos, rápidos e visualmente padronizados, a correção formal passa a funcionar como atalho cognitivo: textos “bonitos”, limpos e sem marcas de hesitação tendem a ser percebidos como mais confiáveis. O efeito é paradoxal. De um lado, a cultura digital fortalece práticas legítimas de variação e criatividade gráfica; de outro, reforça a associação entre ortografia impecável e autoridade discursiva, tensionando a relação entre forma e verdade. É nesse ponto que a discussão sobre massificação da escrita digital se encontra diretamente com o debate sobre pós-verdade e desinformação.

Como mostram \textcite{fazio2015illusorytruth}, já constatamos que a fluência formal aumenta a percepção de verdade, de modo que textos bem escritos tendem a ser julgados como mais confiáveis, mesmo quando falsos. Ensinar ortografia, nesse contexto, é também ensinar desconfiança crítica. Comparar \textit{fake news} mal escritas com outras, impecavelmente ortográficas, por exemplo, ajuda os alunos a perceberem que a correção formal é indício, mas não garante veracidade.

Nesse sentido, analisar a rotulagem algorítmica – como os subtítulos neutros atribuídos pelo Google às figuras historicamente associadas a teorias conspiratórias, conforme apontam \textcite{alrawi2022autocomplete} em seus estudos – torna-se um potente recurso didático para compreender que a legitimidade não está apenas na forma ortográfica, mas também nos mecanismos invisíveis que produzem “credibilidade automatizada”. Investigar com os estudantes por que tais rótulos suavizam discursos extremistas permite discutir vieses, interesses econômicos, processos de curadoria digital e modos de circulação da desinformação. Trata-se de deslocar a ortografia do plano puramente técnico para o campo dos multiletramentos críticos, evidenciando que tanto os textos quanto os resultados de busca precisam ser lidos como construções sociais informadas por algoritmos e relações de poder. Nesse conjunto, a credibilidade deixa de ser atributo puramente formal e se revela como construção sociotécnica, alvo de disputas políticas que exigem leitura crítica dos mecanismos algorítmicos.

A esses pontos soma-se ainda um desafio emergente: a imprecisão dos detectores de IA e a reação que isso provoca entre estudantes. Já há relatos de alunos de graduação que, ao recorrerem a sistemas de escrita assistida, inserem deliberadamente erros ou sintaxe truncada como marca de humanidade, experiências compartilhadas em debates nas disciplinas que ministro, como \textit{Linguística da internet e mídia, escola e letramento digital}. Pesquisas recentes confirmam, ainda, que os detectores apresentam altos índices de falsos positivos, classificando textos humanos como artificiais, sobretudo quando a escrita foge de padrões estatísticos considerados “convencionais”, seja por apresentar frases curtas e repetitivas, vocabulário incomum ou, paradoxalmente, por exibir excesso de fluidez e correção. Nessas situações, como mostram \textcite{saha2025almostai}, até textos humanos apenas revisados ou levemente ajustados podem ser rotulados como totalmente produzidos por IA.

Em minhas próprias experiências, além de já ter tido um texto de minha autoria, integralmente humano e produzido nos anos 2000, rotulado como sendo 100\% produzido por IA, o problema apareceu também no sentido inverso: textos totalmente produzidos por IA foram classificados como apenas parcialmente artificiais, revelando a inconsistência dos critérios utilizados por essas ferramentas \cite{mayers2023aidetectors,saha2025almostai}. Esses desencontros mostram que não se pode tomar a detecção automatizada, nem tampouco a suposta perfeição formal, como garantias absolutas de autoria. Do ponto de vista pedagógico, o erro ganha, então, outro estatuto: em vez de ser apenas falha a corrigir, torna-se espaço não apenas de reflexão, mas de autoria, contraponto necessário à naturalização da escrita “perfeita”, cuja elevação a critério único de qualidade, legitimidade ou humanidade precisa ser problematizada.

Em síntese, ensinar ortografia na era digital implica: (a) trabalhar criticamente com erros e correções automáticas, recuperando seu valor metacognitivo; (b) incorporar práticas digitais como objeto de reflexão escolar, combatendo preconceito linguístico e desvelando exclusões simbólicas; (c) articular ortografia, credibilidade e participação social, formando sujeitos capazes de enfrentar algoritmos, bolhas informacionais e discursos manipuladores; (d) compreender que, em um mundo mediado por tecnologias falíveis de detecção, o erro também pode se transformar em símbolo de autoria, resistência e humanidade.

Dessa forma, os três aspectos delineados – o erro como hipótese, a escrita digital como espaço de variação e a credibilidade como disputa política – configuram um quadro, no qual o ensino da ortografia passa a integrar o campo dos multiletramentos críticos, ampliando sua função de domínio técnico para prática de cidadania discursiva na sociedade digital.


\section{Considerações finais}\label{sec-formato}
Entre algoritmos e símbolos, vimos que a ortografia não é apenas um conjunto de normas gráficas, mas um campo de disputas simbólicas em que se articulam poder, tecnologia e educação. Como destaca \textcite{bourdieu1991language}, a linguagem funciona como capital simbólico, o que, no cenário contemporâneo, marcado pela escrita digital massiva, pela inteligência artificial generativa e pela lógica da pós-verdade, leva a ortografia a se converter em critério de credibilidade e em arena de manipulação discursiva \cite{kalantzis2005learning,fazio2015illusorytruth}. Nesse percurso, mostraram-se decisivos três eixos articulados: a ressignificação do erro, a escrita digital como espaço de variação e disputa, e a credibilidade como construção sociotécnica.

Nesse percurso reflexivo, o erro assume centralidade. Se os algoritmos tendem a apagá-lo antes de emergir, neutralizando sua função epistêmica, a pedagogia crítica pode devolvê-lo ao sujeito como espaço de autoria e reflexão. Em sintonia com a tradição que vai de \textcite{ferreiro1985psicogenese,freire1987pedagogia}, e dialogando com aportes recentes sobre norma digital \cite{fernandezjulia2024normadigital}, o ensino crítico da ortografia passa a disputar com a máquina a possibilidade de ressignificar o erro como experiência de aprendizagem e emancipação. Do mesmo modo, a escrita digital contemporânea, marcada pela circulação híbrida de registros e suportes, evidencia que a ortografia opera também como índice de pertencimento, adequação comunicativa e leitura crítica das práticas sociais.

Este artigo, de caráter teórico-reflexivo, assumiu, assim, como escopo pensar a ortografia no cruzamento entre tecnologias digitais, cultura da desinformação e práticas escolares, apostando em uma abordagem crítica que não se restringe ao domínio técnico, mas que se abre ao horizonte dos multiletramentos \cite{newlondongroup1996multiliteracies,rojo2009letramentos,rojo2012multiletramentos}. Ao trazer a ortografia para o foco desse olhar, enfatizo a necessidade de práticas pedagógicas que considerem a multiplicidade de linguagens e suportes, capazes de formar sujeitos aptos a atuar criticamente em um mundo saturado de algoritmos e discursos manipuladores. Nesse sentido, compreender que a credibilidade textual pode ser estrategicamente produzida – e não apenas derivada da forma – torna-se condição pedagógica para enfrentar a lógica da pós-verdade.

Como desdobramento, torna-se fundamental investigar empiricamente como tais princípios podem ser implementados em situações concretas de ensino, especialmente em contextos que articulam pesquisa e prática docente. Nesse ponto, o Profletras se apresenta como espaço privilegiado: por reunir professores em exercício, por articular formação continuada e investigação aplicada e por incentivar o desenvolvimento de materiais didáticos, projetos de intervenção e análises discursivas. Constitui terreno fértil para avaliar a eficácia de propostas que integrem erro, escrita digital e credibilidade. Futuras pesquisas que mobilizem atividades de sala de aula, objetos digitais de aprendizagem, análises de recepção crítica e experimentações metalinguísticas poderão aprofundar e testar empiricamente as hipóteses aqui delineadas, contribuindo para o avanço do campo.

Em última instância, reafirmo que a ortografia, em tempos digitais, não deve operar como filtro de exclusão, mas como instrumento de inclusão, reflexão e transformação social. Ao articular criticamente erro, variação digital e credibilidade – e ao explorar suas implicações em contextos reais de ensino como os que se constituem no Profletras –, a ortografia se consolida como prática de cidadania discursiva, confirmando seu papel constitutivo na formação de sujeitos críticos e sua relevância para a construção de uma sociedade mais plural e aberta às diferenças.

\printbibliography\label{sec-bib}
% if the text is not in Portuguese, it might be necessary to use the code below instead to print the correct ABNT abbreviations [s.n.], [s.l.]
%\begin{portuguese}
%\printbibliography[title={Bibliography}]
%\end{portuguese}


\begin{dataavailability}
\txtdataavailability{nodata} % options: dataavailable, dataonly, databody, datanotav, nodata
\end{dataavailability}



\end{document}

